\scanpage{203}
\mspage{119}
% archive sheet number: 202
% element: text | top, line 1 | il reste à prouver seulement
il reste à prouver seulement
% element: display | top, line 2 | c prime implies a double prime
\[
c'\Rightarrow a''
\]
% element: text | upper third, lines 3-5 | i.e que F
i.e que $F$ $\pi$-accessible \uline{dans $\widehat C$} implique $\operatorname{card}F(a)\leq\pi\quad\forall a\in C$. Ceci va résulter du résultat un peu plus général

% element: text | upper third, lines 6-8 | Lemme 4.5.8
\uline{Lemme} 4.5.8.\quad $1^\circ$) Dans la catégorie \inserted{$\unsure{\mathscr E|\mathcal E|E}=$} $(\mathrm{Ens})$, les objets $\pi$-accessibles ($\pi$ un cardinal infini) sont les ens. $X$ tels que $\operatorname{card}X\leq\pi$.

% element: text | middle, lines 9-17 | Soit C une catégorie
$2^\circ$) Soit $C$ une catégorie équiv. à une petite catégorie, telle que $\forall x,y\in\mathrm{Ob}\,C$, on ait $\operatorname{card}\mathrm{Hom}(x,y)\leq\pi$. Soit $\mathscr E$ \deleted{\unsure{est|=} $\unsure{\widehat C|C}$}
\deleted{Alors $F$ est $\pi$-accessible, ssi $\forall x\in C$, on \unsure{a} $\operatorname{card}F(x)\leq\pi$}
une $\unsure{\mathcal U|U}$-catégorie \deleted{stable par $\unsure{\varinjlim|\varprojlim}$} \inserted{stable par $\varinjlim$ grandes devant $\pi$, et \unsure{fort|forte}} $\pi$-accessible (i.e. admettant une sous-catégorie génératrice formée d'objets $\pi$-accessibles), \deleted{Alors} \inserted{et stable par produits de cardinal $\leq\pi$}. Alors pour qu'un objet $F$ de $\underline{\mathrm{Hom}}(C,\mathscr E)$ soit $\pi$-accessible \deleted{\unsure{ssi}} \inserted{il faut que} $\forall x\in\mathrm{Ob}\,C$, $F(x)$ soit un objet $\pi$-accessible de $\mathscr E$, et cette condition \uline{est aussi suffisante si} \inserted{$\operatorname{card}\mathrm{Fl}(C)\leq\pi$} \deleted{\unsure{et} \illegible{3} \unsure{les conditions|la condition} $\operatorname{card}\unsure{(\mathrm{Ob}\,C)|(\mathrm{Fl}\,C)}\unsure{\leq\pi|=\pi}$}.
% slanted
% element: margin note | upper left, beside second part of lemma | Cf généralisation
\margin{Cf généralisation 4.14.18 p.\,391}
% slanted
% element: margin note | lower left, linked to sufficient condition | et ce sans la condition
\margin{et ce \uuline{sans} la condition sur $\mathscr E$ pour les produits de cardinal $\leq\pi$.}

% element: text | lower third, lines 18-20 | Dém 1 card X
\uline{Dém}\inserted{(\unsure{suite})}\quad $1^\circ$) $\operatorname{card}X\leq\pi\Rightarrow X$ $\pi$-accessible de $\mathscr E=(\mathrm{Ens})$, car $X=\coprod_{x\in X}\{x\}$, et on applique SGA 4 I 9.9.

% element: text | lower third to foot, lines 21-27 | Inversement
Inversement, supposons $X$ $\pi$-accessible, \deleted{\unsure{on a}} et écrivons $X\simeq\varinjlim_{i\in I}X_i$, où $I$ est l'ens des parties de $X$ qui sont de cardinal $\leq\pi$. Comme $I$ est \unsure{grd|grand} devant $\pi$, l'application identique $X\xrightarrow{\mathrm{id}}X$ se factorise par un des $X_i$, \unsure{donc|d'où} $X_i=X$, donc $\operatorname{card}X=\operatorname{card}X_i\leq\pi$, qed.
